\chapter{Physics as Computation}
\label{sec:physics-as-computation}

Last chapter we went one way:

\begin{center}
\emph{Given WARP + DPO, here’s how physics shows up as invariants and statistics.}
\end{center}

Now we flip it:

\begin{center}
\emph{Given a physical theory, how do we compile it into WARP + DPO so it
\textbf{is} just another region of MRMW?}
\end{center}

This is the "physics compiler" chapter.

We’re not satisfied with saying "physics is computational in spirit."  
We want a literal mapping:

\[
  \text{Physical theory } \mathcal{T}
  \quad\longmapsto\quad
  (\text{WARP } G_\mathcal{T}, \text{ rule set } \mathcal{R}_\mathcal{T})
\]

such that:

- the dynamics of $\mathcal{T}$ show up as rewrite histories of $(G_\mathcal{T},\mathcal{R}_\mathcal{T})$,
- the usual observables of $\mathcal{T}$ show up as queries / coarse-grainings on $G_\mathcal{T}$,
- the symmetries and conservation laws of $\mathcal{T}$ are WARP symmetries and invariants,
- and the whole thing lives in MRMW on equal footing with every other computation.

This is how we stop treating physics as a special external discipline and start
treating it as just another software stack — one with ruthlessly tight tests.

\subsubsection*{26.1  The Compilation Problem}

We’ll treat a (classical) physical theory in a deliberately boring, CS-flavored way:

\begin{itemize}
  \item A \textbf{state space} $S$ (field configurations, particle positions and momenta, etc.).
  \item A \textbf{dynamical law} $F$ (often given as differential equations or an action principle).
  \item A set of \textbf{observables} $\mathcal{O}$ (things agents can measure).
\end{itemize}

The compilation pipeline is:

\[
  (S, F, \mathcal{O})
  \quad\longrightarrow\quad
  (G, \mathcal{R}, C, \mathcal{Q})
\]

where:

\begin{itemize}
  \item $G$ is an initial WARP graph,
  \item $\mathcal{R}$ is a DPO rule set,
  \item $C$ is a coarse-graining / decoding from WARP states to "physical" macrostates,
  \item $\mathcal{Q}$ is a set of query functions implementing observables on $G$.
\end{itemize}

We will call a compilation \textbf{good} if:

\begin{enumerate}
  \item (Fidelity) For relevant time/length scales, $C$ of WARP dynamics matches predictions of $\mathcal{T}$.
  \item (Locality) $\mathcal{R}$ uses bounded-radius patterns; no cheating with global rewrites.
  \item (Symmetry) Key symmetries of $\mathcal{T}$ become graph automorphisms / rule equivariances.
  \item (Sparsity) The WARP representation is not absurdly inflated relative to the physics it encodes.
\end{enumerate}

Bad compilations are allowed — they’re just not interesting. Our goal is to find
and standardize good ones.

\subsubsection*{26.2  World as Graph: Encoding Fields, Particles, Observers}

Step one: turn “the world” into a WARP state.

Three recurring patterns cover most of physics:

\paragraph{(1) Lattice-like fields.}

For continuum-ish theories we can choose a discretization scale and build:

\begin{itemize}
  \item A \textbf{base graph} $G_\text{lat}$ whose nodes represent lattice sites,
        edges represent adjacency (nearest neighbors, etc.).

  \item For each classical field $\phi^a$ (scalar, vector, spin, etc.), we add a
        layer of labels or attached subgraphs representing $\phi^a$'s value at each site.
\end{itemize}

So a scalar field $\phi(x)$ on a 1D lattice becomes:

\begin{itemize}
  \item nodes $i \in \mathbb{Z}$ with edges $(i,i+1)$,
  \item each node $i$ decorated with a tiny graph encoding $\phi_i$ (e.g., a fixed-point
        binary representation, or a symbolic "real" node with approximate arithmetic rules).
\end{itemize}

WARP recursion lets us stack these neatly: "the field" is itself a graph that
knows how to interpret and manipulate per-site subgraphs.

\paragraph{(2) Particle worldlines.}

For particle-centric views, we encode:

\begin{itemize}
  \item \textbf{Particles} as labeled nodes or small complexes (to encode internal degrees of freedom).
  \item \textbf{Worldlines} as sequences of rewrite events on those nodes, or explicit path-like subgraphs.
  \item \textbf{Interactions} as DPO rules that merge/split those structures according to conservation constraints.
\end{itemize}

A scattering process is then literally a local rewrite: a small graph representing
incoming particles is replaced by a small graph representing outgoing particles.

\paragraph{(3) Hybrid: fields on edges, matter on nodes.}

Gauge theories and general relativity-like structures fit more naturally as:

- \textbf{Matter} lives on nodes.
- \textbf{Gauge or geometric fields} live on edges (connections) or higher cells.

So, for a lattice gauge theory:

\begin{itemize}
  \item Nodes: sites with matter fields.
  \item Edges: carry group elements $U_{ij}$ as labels (parallel transport).
  \item Plaquettes: inferred from cycles in the graph; curvature/field strength shows up as holonomy.
\end{itemize}

This is exactly the kind of recursive structure WARP graphs are built for: graphs of graphs of labels.

\paragraph{Observers as subgraphs.}

We won’t pretend observers are mystical. In WARP:

- \textbf{Observers} are just particular subgraphs with stable internal dynamics and:
  \begin{itemize}
    \item sensors (edges from environment into their internal state),
    \item memory (subgraphs that retain structural patterns over time),
    \item and decision rules (rewrite policies over their own subgraph).
  \end{itemize}

They are just as compiled as any other piece of physics; they just happen to ask questions.

\subsubsection*{26.3  From Differential Equations to Rewrite Rules}

Most classical physics is fed to you as a differential equation:

\[
  \frac{\partial \phi}{\partial t} = F\big(\phi, \nabla\phi, \nabla^2\phi, \dots\big)
\]

To compile this into DPO, we do three things:

\begin{enumerate}
  \item \textbf{Discretize spacetime.} Pick a lattice spacing $\Delta x$ and time step $\Delta t$ consistent with stability constraints.
  \item \textbf{Stencil the update.} Express the time derivative as a finite difference update using a local neighborhood.
  \item \textbf{Encode the stencil as a rewrite.} patterns in, patterns out.
\end{enumerate}

Example: 1D wave equation
\[
  \frac{\partial^2 \phi}{\partial t^2} = c^2 \frac{\partial^2 \phi}{\partial x^2}.
\]

Discretized with central differences:

\[
  \phi_i^{(n+1)} = 2\phi_i^{(n)} - \phi_i^{(n-1)} + \lambda^2 \big(\phi_{i+1}^{(n)} - 2\phi_i^{(n)} + \phi_{i-1}^{(n)}\big),
\]
where $\lambda = \frac{c \Delta t}{\Delta x}$.

We turn this into a DPO rule whose left-hand side $L$ matches:

- three time-slices of a site ($\phi_i^{(n-1)}, \phi_i^{(n)}, \phi_i^{(n+1)}$ placeholder),
- plus the neighbors at time $n$ ($\phi_{i-1}^{(n)}, \phi_{i+1}^{(n)}$),

and whose right-hand side $R$ updates the placeholder for $\phi_i^{(n+1)}$ with the computed value.

The "time" here can be:

- an explicit index in the graph (nodes labeled by $(i,n)$),
- or implicit in the sequence of rewrite events (worldline in MRMW).

The scheduler (which sites update when) can itself be encoded as part of the graph:

- synchronous: a global clock subgraph that toggles which slice is being updated,
- asynchronous: local rules that update a site whenever its neighborhood is in a consistent configuration.

Either way, the PDE’s semantics are now just the emergent behavior of repeated local rewrite.

\subsubsection*{26.4  From Action Principles to Local Rewrite}

Many physical theories are cleaner in action form than differential form:

\[
  S[\phi] = \int \mathcal{L}\big(\phi, \partial \phi\big) \, d^dx\, dt,
\]
with dynamics given by $\delta S = 0$.

In WARP-land we:

\begin{enumerate}
  \item \textbf{Discretize the action.} Replace integrals with sums over local motifs (cells, plaquettes, simplices).
  \item \textbf{Attach action contributions} as labels or weights on those motifs.
  \item \textbf{Derive local consistency conditions} that correspond to stationary action.
\end{enumerate}

Concretely:

- Each small subgraph (like a spacetime plaquette) carries a local Lagrangian term $\mathcal{L}_\text{local}$ computed from the fields on its nodes/edges.
- The total discrete action is the sum (or appropriate group combination) of these local contributions.

We can then define rules in two ways:

\begin{description}
  \item[Constraint-based.] Allowed rewrites are those that preserve certain local action conditions (e.g., discrete Euler–Lagrange constraints). Graph states violating them either never occur or get immediately "relaxed" by corrective rules.
  \item[Measure-based.] All rewrites are allowed, but we assign them weights (amplitudes, probabilities) derived from the action difference $\Delta S$. This is the path-integral flavor and ties directly into the "superposition as rewrite bundles" story.
\end{description}

The important point: in action form, you don’t need to guess the local rules by eye. They fall out from asking, "Which local graph changes correspond to a delta in the discrete action that satisfies the extremal principle?"

\subsubsection*{26.5  Worked Example: 1D Wave Equation as WARP graph}

Let’s actually spell one out, even if in sketch.

\paragraph{Encoding.}

\begin{itemize}
  \item Nodes $v_{i,n}$ for each lattice site $i$ and time slice $n$.
  \item Edges between $v_{i,n}$ and $v_{i\pm 1,n}$ (spatial neighbors) and between $v_{i,n}$ and $v_{i,n\pm 1}$ (time neighbors).
  \item Each node $v_{i,n}$ has a subgraph encoding a real value $\phi_{i}^{(n)}$.
\end{itemize}

We maintain a "current slice" marker $T_n$ pointing to the set of nodes at time $n$.

\paragraph{Rule.}

A DPO rule that:

- matches $T_n$, plus a window of nodes $\{v_{i-1,n}, v_{i,n}, v_{i+1,n}, v_{i,n-1}, v_{i,n+1}\}$,
- computes $\phi_i^{(n+1)}$ from the stencil,
- replaces the subgraph at $v_{i,n+1}$ with the updated value,
- and optionally advances the "current slice" marker.

In stringy pseudo-diagram form:

\[
  L: \;\;
  \begin{matrix}
    \phi_{i-1}^{(n)} & \phi_i^{(n)} & \phi_{i+1}^{(n)} \\
    & \phi_i^{(n-1)} & \phi_i^{(n+1)}\_\text{old}
  \end{matrix}
  \quad\Rightarrow\quad
  R: \;\;
  \begin{matrix}
    \phi_{i-1}^{(n)} & \phi_i^{(n)} & \phi_{i+1}^{(n)} \\
    & \phi_i^{(n-1)} & \phi_i^{(n+1)}\_\text{new}
  \end{matrix}
\]

All the usual stability and dispersion discussions are now properties of the rule set:
you can see them as curvature and anisotropy in the induced MRMW neighborhood.

Note what we’ve done: we turned a nice, clean PDE into a pile of graph rewrites, on purpose. Because once it’s in that form, it sits in the same universe as everything else we care about.

\subsubsection*{26.6  Gauge Theories as Graph Rewrite}

Gauge theory is often where "computation" metaphors go to die. Here we keep it simple.

\paragraph{Encoding a lattice gauge theory.}

\begin{itemize}
  \item Nodes: lattice sites with matter fields $\psi_i$.
  \item Edges $(i,j)$: labeled with group elements $U_{ij} \in G$ representing parallel transport.
  \item Plaquettes: cycles of edges; field strength $F$ is encoded as group-valued holonomy around these cycles.
\end{itemize}

\paragraph{Gauge transformation.}

A gauge transformation is a local relabeling:

\[
  \psi_i \mapsto g_i \psi_i, \quad
  U_{ij} \mapsto g_i U_{ij} g_j^{-1}
\]

This is literally a group action on the-labels-of-the-graph. In WARP:

- The map $G \to G$ that applies these relabelings is a symmetry if it commutes with DPO dynamics.
- "Gauge invariance" becomes the requirement that every rule in $\mathcal{R}$ is equivariant under this action.

\paragraph{Dynamics.}

The dynamics (e.g., Wilson action, Kogut–Susskind Hamiltonian) can be compiled by:

- discretizing the action as a sum over plaquettes and sites,
- defining rules that locally adjust $U_{ij}$ and $\psi_i$ to extremize / evolve that action,
- ensuring the rules only depend on gauge-invariant combinations (traces of plaquette holonomies, etc.).

Result: gauge fields are not weird magical objects. They are just labels on edges with group-structured rewrite rules and symmetry constraints.

\subsubsection*{26.7  Quantum Theories and Rewrite Bundles}

We already built the quantum-seeming machinery in Part III: superposition as bundles of rewrite histories, interference as constraint resolution, measurement as minimal-path collapse.

To compile a quantum theory (say, Schrödinger evolution of a wavefunction) we do:

\begin{enumerate}
  \item Encode the classical configuration space (positions, fields) as before.
  \item Introduce a \textbf{bundle structure} over MRMW paths:
    \begin{itemize}
      \item each worldline carries an amplitude,
      \item rules update these amplitudes according to a discretized Hamiltonian or action.
    \end{itemize}
  \item Use a path-integral-style composition:
    \[
      \mathcal{A}(\text{history}) \propto e^{i S[\text{history}]/\hbar}
    \]
    where $S$ is now computed as a functional over rewrite events.
  \item Define observables as queries over the induced probability distribution on coarse-grained states.
\end{enumerate}

In WARP terms:

- The "wavefunction" is not a separate field; it is a weighting over bundles of rewrite paths.
- Schrödinger evolution is just the linear update of these weights according to local rules derived from the Hamiltonian.
- Interference is what you get when different rewrite histories end in the same coarse-grained graph and their amplitudes combine.

This makes quantum vs classical a choice of \emph{how} you traverse MRMW, not \emph{what} MRMW is.

\subsubsection*{26.8  Observables as Graph Queries}

We want a clean dictionary between physics lab talk and WARP ops.

Given a compiled pair $(G,\mathcal{R})$:

\begin{itemize}
  \item A \textbf{physical observable} $O$ is a function on (coarse-grained) graphs:
    \[
      O : \text{WARP states} \to \mathbb{R}^k \text{ or some measurable space}.
    \]
  \item A \textbf{measurement procedure} is a composed process:
    \begin{enumerate}
      \item Evolve $(G,\mathcal{R})$ for some time.
      \item Apply a coarse-graining $C$.
      \item Apply $O$ to the result.
    \end{enumerate}
\end{itemize}

This sounds trivial until you weaponize it:

- We can distinguish observables that are \emph{local} queries (depend on limited subgraphs) from \emph{global} ones.
- We can explicitly see which observables commute (their queries touch disjoint or causally separated regions) and which don’t.
- We can model "apparatus" not as magical projection operators, but as additional WARP graph subgraphs that implement $C$ via interaction.

In other words, the measurement problem becomes a design problem: what queries do your observer-subgraphs implement, and how do they distort the underlying dynamics?

\subsubsection*{26.9  What Makes a Compilation "Good"?}

Just because you can encode a theory as WARP + DPO doesn’t mean you did it well.
Good compilations have:

\paragraph{(1) Structural locality.}

Rules have small patterns and bounded match radius. This:

- ensures a causal graph with meaningful light cones,
- keeps implementation cost reasonable,
- matches the "no action at a distance" vibe of real physics.

\paragraph{(2) Symmetry transparency.}

Symmetries of the physical theory show up as genuine graph automorphisms / rule
equivariance, not as accidental numerical coincidences.

If a compilation hides symmetry in boundary conditions or gross hacks, it’s suspect.

\paragraph{(3) Coarse-graining friendliness.}

There exists a natural $C$ such that:

- entropy and thermodynamic behavior are visible,
- macroscopic variables behave smoothly,
- the arrow of time shows up in the expected direction.

If everything is brittle under coarse-graining, you haven’t captured the right structures.

\paragraph{(4) Rulial robustness.}

Small changes in rules should correspond to physically interpretable perturbations
(e.g., slightly different coupling constants, different fields), not total nonsense.

That is: nearby points in rulial space correspond to nearby physical theories.
If not, you’ve chosen a pathological coordinate chart on MRMW.

These criteria are how CΩMPUTER decides which physical encodings are worth caring about.

\subsubsection*{26.10  This is Not "Just Simulation"}

Let’s be explicit about the difference.

\paragraph{Simulation mindset.}

- Pick a theory.
- Discretize it.
- Write a program that approximates its evolution.
- Use the program as a black-box calculator.

\paragraph{CΩMPUTER mindset.}

- Treat that discretization as a first-class \emph{universe} in MRMW.
- Put it on the same footing as databases, neural networks, protocols, etc.
- Compare it geometrically to other universes: what’s near it? what’s equivalent to it?
- Build machines (from Part IV) that span across this universe and others.

The point is not "we can run physics on a computer." That’s old news.
The point is:

\begin{quote}
\textbf{Physics is now an API over MRMW.}
\end{quote}

Once a physical theory is compiled into WARP + DPO:

- you can compose it with other universes (e.g., a market model, a biological model),
- you can apply rulial analysis (how does complexity scale? where does curvature spike?),
- you can subject it to adversarial universes, counterfactual engines, and optimization across worlds.

You stop drawing a sacred line between "the real world" and "the models".
It’s universes all the way down.

\subsubsection*{26.11  On to CΩSMOS}

With this chapter, we now have:

\begin{itemize}
  \item a way to read physics out of computation (previous chapter),
  \item and a way to push physics back into computation (this chapter).
\end{itemize}

The next step is obvious and a little insane:

\begin{center}
\emph{Apply this compiler to the one actual universe we live in.}
\end{center}

\noindent
That’s what the next chapter, \emph{CΩSMOS: The Physical Universe as WARP graph},
attempts: not a precise final answer, but a concrete sketch of how a realistic
candidate $(G,\mathcal{R})$ for our cosmos might look, and what it would mean
to treat "the universe" as just one WARP instance running inside CΩMPUTER.
